\section{Scope of the Finite-Variance Hypotheses}
\label{sec:scope-of-the-finite-variance-hypotheses}

The zero-set row-condition classification of
\cref{sec:zero-set-row-condition-and-linear-equivalence} shows that
positivity is not merely a regularity convenience: without it,
Hilbert--Schmidt proximity must be paired with the zero-set row condition,
and boundary geometry may collapse the equivalence group to permutations.
The remaining assumptions control the Fourier/Fisher mechanism for bounded
invertible linear mixing.

The finite-variance a.e.-positive-density theorem is broad within its regime.
It includes Laplace densities centered with variance one despite their
cusp at the origin, as well as densities centered with variance one and proportional to
\(e^{-x^4}\).
Symmetry is not needed in any step of the proof.

For examples allowing both asymmetry and heavy tails, take a density
proportional to
\begin{equation}
 (1+y^2)^{-(\tau+1)/2}
                           e^{\vartheta\arctan y},
 \qquad \tau>2,\quad \vartheta\in\R,
 \label{eq:asymmetric-heavy-tail-example}
\end{equation}
and center and rescale to variance one.
Its location score before centering and rescaling is
\[
 \frac{(\tau+1)y-\vartheta}{1+y^2}.
\]
Both this score and its product with \(y\) are bounded.
The density is everywhere positive and smooth.
It has a finite second moment.
Centering and rescaling preserve finiteness of the location and scale
Fisher information.
So the density centered with variance one satisfies \eqref{eq:main-assumptions}.
For \(2<\tau\le4\), its fourth moment is infinite.
When \(\vartheta\ne0\), the two tails have unequal leading constants.
So centering does not make the density even.
The case \(\vartheta=0\) is a rescaled Student \(t\) family.
These examples show that neither evenness nor a fourth moment is implicit in
the theorem.

The Gamma example used to illustrate the rotational Fourier defect in
\cref{sec:variance-and-affine-normalization} should therefore be read
through the zero-set row-condition theorem, not through the
a.e.-positive-density theorem.
Its centered density vanishes on a half-line.
The rotational Fourier defect still gives a local Hellinger obstruction to rotations,
but equivalence is decided by the half-line support classification:
homogeneous linear equivalence permits only coordinate permutations.
The same support theorem classifies all bounded coordinate laws, including
disconnected supports and laws without densities.
\Cref{cor:affine-gamma-classification} treats the distinct affine Gamma problem,
where support-preserving scales do allow equivalence.

The moment assumptions play a different role from the boundary assumptions.
Support controls which rows are admissible from the viewpoint of null sets;
the second moment controls whether arbitrary \(\ell^2\) rows define the
measurable linear action used throughout \cref{thm:measure-equivalence}.
It also has a specific role beyond differentiating the expectation of the
product of centered complex exponentials.
It makes every \(\ell^2\) row into a convergent random series.
It supplies uniform integrability of squares of linear forms.
It gives the coefficient-subsequence argument with an independent Gaussian
summand in \cref{lem:unit-row-rigidity}.
Thus the first-moment ODE result is stronger in its moment scope than
\cref{thm:measure-equivalence}.

Under \eqref{eq:main-assumptions}, the Hellinger path estimates give
\(L^1\) convergence and uniform integrability of likelihood ratios.
The finite real Fourier sums are bounded in each coordinate.
Their dimension-free bounds are the \(L^2(\mu_n)\) and generator estimates
proved above.
No variance bound for their aggregate under an arbitrary mixed output law is
needed in the tail-symmetry argument.
